\SourceLeaf{291}%
% A544-P0291-D01
ellipsin, cuius aequatio est
\[
aggpp+2bghpq+chhqq=1
\]
iacentium semper adnumeratis si quae sunt in peripheria ipsa. Cuius ellipsis area $=\frac\pi{gh\surd(ac-bb)}=\frac\pi{gh\surd D}$ erit limes, ad quem productum $m'\lambda\lambda=\frac{m'}M$ in infinitum appropinquabit, decrescente $\lambda$ vel crescente $M$ in infinitum.

% A544-P0291-D02
Ceterum manifestum est, theorema nostrum complecti casum ubi alterutra indeterminatarum $x,y$ sola per saltus progredi debet, dum alterius valor nulli conditioni subiicietur. Patet enim, hoc idem esse, ac si vel $h$ vel $g$ statuatur $=1$.

% A544-P0291-D03
\Article{7}
Quae hactenus exposita sunt, ab indole coëfficientium formae $axx+2bxy+cyy$ sunt independentia: abhinc vero supponemus, hosce coëfficientes esse integros. Ita quaevis combinatio valorum integrorum quantitatum $x,y$ ipsi formae valorem integrum conciliabit, sive repraesentationi alicuius numeri integri per istam formam respondebit. Hinc patet, complexum omnium combinationum valorum integrorum quantitatum $x,y$, per quos forma $F=axx+2bxy+cyy$ valorem non maiorem limite $M$ nanciscatur, esse idem ac complexum omnium repraesentationum numerorum integrorum limitem $M$ non egredientium, sive usque ad hunc limitem incl., si ipse est numerus integer. Quodsi itaque brevitatis gratia multitudinem repraesentationum diversarum numeri determinati integri $n$ per formam $F$ per $F(n)$, vel quatenus ambiguitas non metuenda simpliciter per $Fn$ denotamus, numerus supra per $m$ expressus erit $=F0+F1+F2+F3+\text{etc.}+FM$, theoremaque primum sequentem induit formam.

% A544-P0291-D04
\textsc{Theorema III.}\quad\emph{Aggregatum $F0+F1+F2+\text{etc.}+FM$ proxime exprimitur per $\frac{\pi M}{\surd D}$, approximatione in infinitum crescente, dum $M$ in infinitum augetur.}

% A544-P0291-D05
\Article{8}
Theoremati tertio repraesentationes \emph{omnium} numerorum spectanti aliud adiungere convenit, solos numeros impares spectans. Manifesto per formam $F$ numeri impares repraesentari nequeunt, si $a$ et $c$ simul sunt numeri pares: quapropter disquisitio ad tres reliquos casus restricta erit.

% A544-P0291-D06
I. Quoties $a$ est impar, $c$ par, numerus impar repraesentatur, tribuendo ipsi $x$ valorem imparem, valore ipsius $y$ arbitrario manente. Theorema II. 
\SourceLeaf{292}%
% A544-P0292-D01
itaque, statuendo $g=2,G=1,h=1$, docet, multitudinem omnium combinationum valorum talium ipsorum $x,y$, qui formae valorem imparem limite $M$ non maiorem concilient, approximatione infinita exprimi per $\frac{\pi M}{2\surd D}$, crescente $M$ in infinitum.

% A544-P0292-D02
II. Quoties $a$ est par, $c$ impar, ad repraesentationem numeri imparis requiritur, ut $y$ sit impar, unde statuendo $g=1,h=2,H=1$, ad eandem conclusionem deferimur.

% A544-P0292-D03
III. Quoties tum $a$ tum $c$ impar est, vel valor impar ipsius $x$ cum valore pari ipsius $y$ combinari debet, vel valor par ipsius $x$ cum valore impari ipsius $y$, ut prodeat valor impar formulae. Multitudo omnium combinationum tum prioris generis tum posterioris, pro quibus valor formae limitem $M$ non egreditur, approximatione infinita per $\frac{\pi M}{4\surd D}$ exprimitur, quapropter multitudo omnium combinationum, quae formae valores impares limitem $M$ non egredientes producunt, etiam hic approximatione infinita per $\frac{\pi M}{2\surd D}$ exprimitur.

% A544-P0292-D04
Iam quum complexus omnium talium combinationum nihil aliud sit, nisi complexus omnium repraesentationum omnium numerorum $1,3,5,7\ldots M$, quoties $M$ est integer impar, vel $1,3,5,7\dotsfour M-1$, quoties $M$ est par, habemus

% A544-P0292-D05
\textsc{Theorema IV.}\quad\emph{Aggregatum}
\[
F1+F3+F5+F7\ldotp\ldotp+FM\quad\text{vel}\quad F1+F3+F5+F7\ldotp\ldotp+F(M-1)
\]
\emph{(prout $M$ impar est vel par) approximatione infinita exprimitur per $\frac{\pi M}{2\surd D}$, siquidem $F$ est forma, in qua alteruter coëfficientium $a,c$ vel uterque est impar.}
\BookRule{.30\textwidth}

% A544-P0292-D06
\par\bigskip\pdfbookmark[0]{De nexu. Fragmentum III}{nexus-fragment-III}\centerline{[III.]}\medskip
Es sei $C$ der Complexus der Repräsentanten sämmtlicher Classen der formae proprie primitivae für den Determinant $-D$. Wir bezeichnen durch $(n)$ die Anzahl aller Darstellungen der Zahl $n$ durch Formen aus dem Complexus $C$. Es sei $p$ eine ungerade Primzahl. Dann ist

\SourceLeaf{293}%
% A544-P0293-D01
\[
\begin{aligned}
&1.\quad(pn)=(n)\quad\text{wenn }p\text{ ein Divisor von }D\\
&\left.\begin{aligned}2.\quad(pn)&=(n)+(h)\\3.\quad(pn)&=-(n)+(h)\end{aligned}\right\}
\ \text{wenn }p\text{ Nichtdivisor von }D\\
&\hspace{7em}\left\{\begin{aligned}&\text{Divisor von }xx+D\\&\text{Nichtdivisor von }xx+D\end{aligned}\right.
\end{aligned}
\]
wo $n=hp^\mu$, $\mu$ beliebig und $h$ nicht durch $p$ theilbar.
\[
\begin{aligned}
\text{Im Fall }1.\ \text{ist}\quad&(h)=(ph)=(pph)=(p^3h)\text{ etc.}\\
2.\quad&(ph)=2(h),\ (pph)=3(h),\ (p^3h)=4(h)\text{ etc.}\\
3.\quad&(ph)=0,\ (pph)=(h),\ (p^3h)=0,\ (p^4h)=(h)\text{ etc.}
\end{aligned}
\]
\begin{center}$\ast\qquad\ast$\par$\ast$\end{center}

% A544-P0293-D02
Aus jeder Classis pr. pr. für den Determinans $=-D$, deren Anzahl $=\lambda$, sei eine Form ausgewählt, und der Complexus dieser Form sei $L$.

% A544-P0293-D03
Man bezeichne durch $fA$ die Anzahl sämmtlicher Darstellungen der Zahl $A$ durch Formen aus $L$.

% A544-P0293-D04
Es sei ferner $f(A;p)=f\frac A{p^\alpha}$, wenn $p^\alpha$ die höchste Potenz der Primzahl $p$ ist, welche $A$ misst; ferner $f(A;p,q)=f\frac A{p^\alpha q^\beta}$, wenn $q$ eine andere Primzahl deren höchste $A$ messende Potenz $=q^\beta$ und so ferner $f(A;p,q,r)=f\frac A{p^\alpha q^\beta r^\gamma}$ wenn $r$ eine dritte Primzahl, deren höchste Potenz $A$ messend $r^\gamma$ ist u. s. w.
\BookRule{.30\textwidth}

% A544-P0293-D05
\par\bigskip\pdfbookmark[0]{De nexu. Fragmentum IV}{nexus-fragment-IV}\centerline{[IV.]}\medskip
Man bezeichne durch $(n)$ die Anzahl der Werthe $x$ aus dem Complexus
\[
0,1,2,3,4\dotsfour p^n-1
\]
für welche $xx-D=xx-ap^\mu$ durch $p^n$ theilbar ist.
\[
\begin{array}{c@{\qquad}cc}
1)\ \mu\text{ ungerade z. B. }=7.&\multicolumn2c{2)\ \mu\text{ gerade z. B. }=6}\\
(1)=1&aNp&aRp\\
(2)=p&(1)=1&(1)=1\\
(3)=p&(2)=p&(2)=p\\
(4)=pp&(3)=p&(3)=p\\
(5)=pp&(4)=pp&(4)=pp
\end{array}
\]

\SourceLeaf{294}%
% A544-P0294-D01
\[
\begin{array}{c@{\qquad}c@{\qquad}c}
(6)=p^3&(5)=pp&(5)=pp\\
(7)=p^3&(6)=p^3&(6)=p^3\\
(8)=0&(7)=0&(7)=2p^3\\
(9)=0&(8)=0&(8)=2p^3\\
\text{etc.}&\text{etc.}&\text{etc.}
\end{array}
\]

% A544-P0294-D02
\[
\begin{array}{c@{\qquad}c}
\text{Man mache nun}&\text{Dann ist}\\[.4em]
(1)-\dfrac{(2)}p=(1)'&fp=(1)'\\[.4em]
(2)-\dfrac{(3)}p=(2)'&fpp=1+(2)'\\[.4em]
(3)-\dfrac{(4)}p=(3)'&fp^3=(1)'+(3)'\\[.4em]
(4)-\dfrac{(5)}p=(4)'&fp^4=1+(2)'+(4)'\\
\text{etc.}&\text{etc.}
\end{array}
\]
Es ist folglich, $\frac{p-1}p(1+\frac{fp}p+\frac{fpp}{pp}+\frac{fp^3}{p^3}+\text{etc.})=T$ gesetzt,
\[
\begin{aligned}
\frac{p+1}p T&=1+\frac{(1)'}p+\frac{(2)'}{pp}+\frac{(3)'}{p^3}+\frac{(4)'}{p^4}+\text{etc.}\\
&=1+\frac{(1)}p=1+\frac1p
\end{aligned}
\]
Also $T=1$
\BookRule{.30\textwidth}

% A544-P0294-D03
\par\bigskip\pdfbookmark[0]{De nexu. Fragmentum V}{nexus-fragment-V}\centerline{[V.]}\medskip
Multitudo classium mediocris\begingroup\renewcommand{\thefootnote}{*)}\footnote{[Vergl. \emph{Disquiss. Arithm.} art. 302; die dortige Formel weicht um eine Constante $\delta$ von der hier im Text vorkommenden ab.]}\endgroup\ circa determinantem negativum $-D$ est proxime
\[
=\frac{\pi\surd D}{4(1+\tfrac1{27}+\tfrac1{125}+\text{etc.})}
\]
Multitudo vera exprimitur sequentibus formulis, ubi brevitatis caussa scribitur $m$ pro multitudine mediocri, $M$ pro vera; $p,q$ exprimunt omnes numeros impares primos ipsum $D$ non metientes, ille divisores, hic non-divisores ipsius $\square+D$; $r$ numeros\begingroup\renewcommand{\thefootnote}{**)}\footnote{[impares.]}\endgroup\ primos ipsum $D$ metientes:

\SourceLeaf{295}%
% A544-P0295-D01
\[
\begin{aligned}
\text{I.}\quad M&=m\ \text{Prod. ex }\frac{p^3+p^2}{p^3-1}\pdot\frac{q^3-q^2}{q^3-1}\pdot\frac{r^3-r}{r^3-1}\\
\text{II.}\quad M&=\frac{\pi\surd D}4\ \text{Prod. ex }\frac{p+1}p\pdot\frac{q-1}q\pdot\frac{rr-1}{rr}\\
\text{III. NB.}\quad M&=\frac{2\surd D}\pi\ \text{Prod. ex }\frac p{p-1}\pdot\frac q{q+1}\\
\text{IV.}\quad M&=\surd\left\{\frac D2\pdot\text{Prod. ex }\frac{p+1}{p-1}\pdot\frac{q-1}{q+1}\pdot\frac{rr-1}{rr}\right\}\\
\text{V.}\quad M&=\frac{2\surd D}\pi\{1\pm\tfrac13\pm\tfrac15\pm\tfrac17\pm\tfrac19\text{ etc.}\}
\end{aligned}
\]
NB. Die Formel III wird unmittelbar aus der Vergleichung der beiden Arten, die darstellbaren Zahlen bis zu einer gewissen Grenze zu zählen, abgeleitet.
\BookRule{.30\textwidth}

% A544-P0295-D02
\par\bigskip\pdfbookmark[0]{De nexu. Fragmentum VI}{nexus-fragment-VI}\centerline{[VI.]}\medskip
\textsc{Theorema.} Multitudo classium, in quas omnes formae binariae proprie primitivae determinantis negativi $-D$\begingroup\renewcommand{\thefootnote}{*)}\footnote{[distribuuntur.]}\endgroup, aequalis est
\[
\frac\pi4\times\surd D\times\text{Prod. ex. }\frac{p-1}p\frac{q+1}q\times\frac{rr-1}{rr}
\]
designantibus
\begin{description}
\item[$p$] omnes numeros primos\begingroup\renewcommand{\thefootnote}{**)}\footnote{[impares.]}\endgroup\ quorum non-res. est $-D$
\item[$q$] omnes numeros primos\textsuperscript{**)} quorum res. $-D$
\item[$r$] omnes numeros primos\textsuperscript{**)} ipsum $D$ metientes
\end{description}
\[
=\frac{\frac\pi4\surd D\ \text{Prod. ex }\frac{rr-1}{rr}}{1\pm\tfrac13\pm\tfrac15\text{ etc.}}
\]

% A544-P0295-D03
ubi in denom. signum posit. praeponitur fractt., quarum denom. sunt in forma non divis.; negat. iis, quarum denom. sunt in forma divisorum ipsius $xx+D$; eae vero, quarum denom. ad $D$ non forent primi, omnino omittuntur\begingroup\renewcommand{\thefootnote}{***)}\footnote{[Bezeichnet man mit $m$ alle positiven ganzen Zahlen, die relative Primzahlen zu $2D$ sind, und benutzt man das durch \textsc{Jacobi} verallgemeinerte Symbol von \textsc{Legendre}, so ist die obige Regel für die Zeichenbestimmung in folgender Weise zu berichtigen: in der vorhergehenden Formel ist der Nenner
\[
1\pm\tfrac13\pm\tfrac15\text{ etc.}=\Sigma\pm\left(\frac{-D}m\right)\frac1m
\]
wo das obere oder untere Zeichen zu nehmen ist, je nachdem die Zahl $m$ ein Product aus einer geraden oder ungeraden Anzahl (gleicher oder ungleicher) Primzahlen ist; dagegen ist im Zähler der nachfolgenden Formel
\[
1\pm\tfrac13\pm\tfrac15\ldots=\Sigma\left(\frac{-D}m\right)\frac1m
\]
]}\endgroup.

\SourceLeaf{296}%
% A544-P0296-D01
\[
\begin{aligned}
&=\frac{2\surd D(1\pm\tfrac13\pm\tfrac15\ldots)}\pi\\
&=\frac{\operatorname{cotg}\theta\pm\operatorname{cotg}3\theta\pm\operatorname{cotg}5\theta\ldots\pm\operatorname{cotg}n\theta}{N:\surd D}
\end{aligned}
\]
ponendo $\theta=\frac\pi N$, $N=\left\{\begin{smallmatrix}1\\4\end{smallmatrix}\right\}D$ et ponendo pro $n$ omnes numeros ad $D$ primos signo ut supra determinato\begingroup\renewcommand{\thefootnote}{*)}\footnote{[Siehe die weiter unten folgende Note zu diesem Fragment.]}\endgroup.

% A544-P0296-D02
Pro determ. pos. erit mult. Classium\begingroup\renewcommand{\thefootnote}{**)}\footnote{[In der nachfolgenden Formel bedeutet $D$ den positiven Determinanten, und es ist
\[
1\pm\tfrac13\pm\tfrac15\ldots=\Sigma\left(\frac Dm\right)\frac1m
\]
]}\endgroup
\[
=\frac{2\surd D(1\pm\tfrac13\pm\tfrac15\ldots)}{\log .T+U\surd D}
\]
Designantibus $T,U$ valores minimos quantitatum $t,u$ aequationi $tt-Duu=1$ satisfacientes
\[
=\frac{\log\sin\tfrac12\theta\pm\log\sin\tfrac32\theta\pm\log\sin\tfrac52\theta\text{ etc.}}{\log T+U\surd D}
\]
\BookRule{.30\textwidth}

% A544-P0296-D03
\par\bigskip\pdfbookmark[0]{De nexu. Fragmentum VII}{nexus-fragment-VII}\centerline{[VII.]}\medskip
Pro determinante negativo $-p$, qui\begingroup\renewcommand{\thefootnote}{***)}\footnote{[d. h. wenn $p$ eine positive Primzahl von der Form $4n+1$ ist.]}\endgroup\ est numerus primus formae $4n+1$, multitudo classium est\begingroup\renewcommand{\thefootnote}{$\dagger$)}\footnote{[multitudo classium est $=2(\alpha-\beta)$.]}\endgroup\ $=(\alpha-\beta)$, ubi $\alpha$ multitudo residuorum quadraticorum in quadrante primo
\[
1\pdot2\pdot3\dotsfour\tfrac14(p-1)
\]
$\beta$ multitudo non-residuorum.
\BookRule{.30\textwidth}

